% 生成AI等を利用していない場合は，このファイルをコメントだけのままにします．
% その場合，見出しも空白頁も出力されません．
%
% 生成AI等を利用した場合は，提出時点の大学・研究科・専攻の規則と
% 指導教員の指示を確認したうえで，次の雛形をコメント解除し，
% 角括弧内を実際の利用内容へ置き換えてください．利用していない用途や，
% 自分で確認していない事項を記載してはいけません．
%
% \chapter*{生成AI等の利用に関する申告}
% \markboth{生成AI等の利用に関する申告}{生成AI等の利用に関する申告}
% \addcontentsline{toc}{chapter}{生成AI等の利用に関する申告}
%
% 本論文の作成にあたり，著者は［サービス名，モデル名，版および利用期間］を，
% ［論点整理，プログラム実装の補助，コードの整理，文章構成，表現の推敲，
% 批判的確認等の実際の用途］に使用した．入力した情報は
% ［入力情報の範囲と，外部サービスへ入力しなかった情報］である．
% 生成AI等の出力はそのまま研究上の根拠として採用せず，
% ［実験設計，ソースコード，実行結果，数値，図表，主張，引用文献，
% 最終稿等の実際に確認した対象］を，著者が
% ［原資料，一次文献，実行ログ，再実行等の実際の方法］により確認した．
% 生成AI等を著者として扱わず，本論文の正確性，研究公正，独創性および
% 最終内容については著者が責任を負う．
%
% 生成AI等がデータ生成，前処理，解析，推論，評価者または判定器等として
% 研究方法そのものに関与した場合は，この申告だけで済ませず，
% 方法・実験章にもモデル，版，入力，出力，設定，検証方法および限界を
% 再現可能な範囲で記載してください．
